\chapter*{Conclusion générale}
\addcontentsline{toc}{chapter}{Conclusion générale}
\markboth{\textbf{Conclusion générale}}{}

Ce travail avait pour ambition de concevoir un système capable de produire, en temps réel et sur le lieu de vente, des recommandations commerciales et des décisions de réapprovisionnement justes, contextualisées, explicables et contrôlées. Le cycle CRISP-DM a été parcouru dans son intégralité, chaque phase close par une décision explicite autorisant la suivante.

\section*{Bilan du travail réalisé}

La \textbf{compréhension métier} a produit sept objectifs métier traduits en huit objectifs analytiques, chacun assorti d'une métrique et d'un critère de succès fixé avant toute modélisation. Cette antériorité, tenue sans exception, est ce qui rend l'évaluation du chapitre~\ref{chap:evaluation} crédible.

La \textbf{compréhension et la préparation des données} ont réuni quatre familles de sources en une base unique de cinquante et une tables, portant près de deux millions de transactions sur vingt-deux mois. L'exploration a joué son rôle méthodologique en validant certaines hypothèses et en invalidant d'autres : la découverte que 572 références seulement, sur les 4 593 du catalogue, ont fait l'objet d'une vente, et que 1,7 \% d'entre elles portent 80 \% du volume, a conduit à resserrer explicitement le périmètre de la prévision de demande. Neuf anomalies de qualité ont été relevées et traitées, ou assumées lorsqu'elles ne pouvaient pas l'être.

La \textbf{modélisation} a construit un moteur de séries temporelles déterministe produisant son diagnostic en moins d'une seconde, un modèle appris global organisé en cascade de repli, une prévision de demande en deux étages séparant signaux lents et signaux rapides, et des modèles de décision de stock entièrement déterministes. Huit agents intelligents ont été conçus, coordonnés par un graphe d'états à mémoire partagée dont la discipline d'écriture rend le parallélisme possible. Un contrôle de conformité à sept règles et un point d'arrêt humain obligatoire encadrent toute production du système.

L'\textbf{évaluation} a établi que le modèle appris global réduit de 27,9 \% l'erreur du moteur statistique sur un rétro-test de 24 681 prévisions réparties en six plis chronologiques, en battant l'alternative sur 150 des 151 points de vente, avec un biais quasi nul. La classification du risque de stock est cohérente sans exception. Le contrôle de conformité présente une exactitude parfaite et, surtout, aucun faux blocage.

Le \textbf{déploiement} a livré un système en fonctionnement, dont la chaîne causale est fermée de la vente qui épuise le stock jusqu'à la réception de la commande qui le reconstitue.

\section*{Apports}

Trois apports méritent d'être soulignés.

Le premier est le \textbf{principe de séparation entre calcul et formulation}. Aucune valeur engageant une décision — quantité, seuil, score, verdict — n'est produite par un modèle génératif : cinq des neuf unités du système n'invoquent aucun modèle de langage. Cette propriété n'est pas une précaution de style ; elle est ce qui rend la totalité des calculs décisionnels couverte par des tests unitaires reproductibles, et donc ce qui rend une recommandation auditable par le responsable qui doit l'assumer.

Le deuxième est le \textbf{couplage effectif des deux domaines métier}. L'exploration a montré que disponibilité et potentiel commercial sont largement décorrélés : c'est précisément cette absence de redondance qui donne sa valeur au score multicritère, et qui permet au système de ne jamais recommander un produit indisponible.

Le troisième est le \textbf{dispositif de mesure en exploitation}. Un système d'intelligence artificielle dont personne ne mesure la qualité perçue dérive sans que quiconque s'en aperçoive. Le dispositif mis en place a produit le résultat le plus instructif de ce travail, exposé ci-dessous.

\section*{Limites}

Ce travail présente des limites qu'il serait malhonnête de taire.

La plus significative est l'\textbf{écart entre performance de banc et performance réelle}. La qualité mesurée en exploitation s'établit à 3,38 sur 5, contre 4,76 sur les cas de banc du même domaine — plus d'un point d'écart. Les cas de banc sont propres, les cas réels sont sales : contexte partiel, situations ambiguës, références mal paramétrées. La conséquence méthodologique dépasse ce projet : \textbf{un banc hors ligne mesure une borne supérieure, pas une performance}.

La deuxième est le \textbf{taux d'acceptation des propositions}, à 48,1 \% sur quatre-vingt-une décisions tracées. Le critère fixé au chapitre~\ref{chap:comprehension-metier} exigeait une majorité ; il n'est pas atteint. L'analyse différenciée montre que le déficit se concentre sur les stratégies commerciales et non sur les décisions de stock, ce qui oriente précisément la correction. Le système fonctionne sans encore convaincre — et cette distinction, que seule l'évaluation métier permet d'établir, est l'enseignement principal du chapitre~\ref{chap:evaluation}.

La troisième est la \textbf{validation incomplète du juge automatique}. Les trois contrôles prévus — discrimination sur gabarit, détection de chiffre corrompu, déterminisme — n'ont pas été exécutés. En conséquence, toutes les mesures qualitatives de ce rapport doivent être lues comme indicatives et non comme établies. Cette lacune est reconnue dans la revue de processus.

La quatrième est la \textbf{qualité du paramétrage logistique}. Les délais fournisseurs, coûts de possession et de passation sont, pour une large part du catalogue, laissés à leur valeur par défaut. La méthode de calcul est vérifiée ; les valeurs produites restent des ordres de grandeur défendables, non des optima.

Enfin, une observation transversale mérite d'être rapportée : sur les neuf difficultés majeures rencontrées, \textbf{sept relèvent de la donnée et de l'infrastructure, deux seulement de la modélisation}. Cette répartition contredit l'intuition selon laquelle la difficulté d'un projet d'intelligence artificielle réside dans le choix des algorithmes.

\section*{Perspectives}

Les perspectives sont hiérarchisées par leur effet attendu sur le seul critère métier non atteint.

En \textbf{priorité haute} : conduire la validation du juge automatique, sans laquelle aucune mesure qualitative n'est probante ; enrichir le corpus documentaire, dont la faible couverture limite simultanément la pertinence du coaching et la qualité de la génération ancrée — c'est le levier unique le plus efficace sur le taux d'acceptation ; renseigner les paramètres logistiques, condition pour que les quantités recommandées cessent d'être des ordres de grandeur.

En \textbf{priorité moyenne} : corriger l'incohérence de consolidation de l'indicateur de chiffre d'affaires du jour ; ramener le 95\textsuperscript{e} centile de la latence conversationnelle dans une plage compatible avec l'usage au comptoir ; rejouer le rétro-test avec réentraînement par pli afin de fixer la fréquence de réentraînement en exploitation.

En \textbf{priorité basse} : étendre le périmètre à l'ensemble du réseau, ce que l'architecture permet sans modification ; ajouter les fonctions classées \textit{Won't have} dans le backlog initial.

\section*{Enseignement méthodologique}

Conduire ce projet par CRISP-DM plutôt que de simplement citer la méthode a eu un effet concret et vérifiable : c'est parce que les critères de succès avaient été fixés avant la modélisation qu'il a fallu reconnaître, en fin de parcours, que deux d'entre eux n'étaient pas atteints. Un rapport d'évaluation dont tous les seuils sont franchis est plus souvent le signe de seuils ajustés après coup que celui d'un système exceptionnel. La discipline imposée par la méthode aura donc principalement servi à rendre le travail honnête — ce qui est, pour un projet destiné à être repris, sa qualité la plus utile.
